\section{Model, network class, and main results}\label{sec:model-results}

\subsection{Scalar leaky-recovery RFDE}

Let \(r=5\sqrt5\), \((\tau_0,\tau_1)=(4\sqrt5,5\sqrt5)\), and
\(H(v)=(v-1)^3\).  The autonomous scalar model is
\begin{align}
 \dot v={}&v-\frac{v^3}{3}-w
 +\eps\kappa_1\left\{\frac{v(t-\tau_0)+v(t-\tau_1)}2-v(t)\right\}
 \notag\\
 &+\eps\kappa_3\left\{
 \frac{H(v(t-\tau_0))+H(v(t-\tau_1))}{2}-H(v(t))\right\},\label{eq:scalar-v}\\
 \dot w={}&\eps(v-a-w).\label{eq:scalar-w}
\end{align}
The validated center is
\begin{equation}
 \eps=\frac15,\qquad a=\frac14,\qquad
 \kappa_1=\frac1{250},\qquad \kappa_3=\frac1{200}.       \label{eq:center}
\end{equation}
A physical pulse of amplitude \(J\) adds
\(J\mathbf1_{[0,1)}(t)\) to \cref{eq:scalar-v}; after release the vector
field is exactly autonomous again.

The full phase space is
\(X=C([-r,0],\R^2)\).  Because the recovery equation has no delayed recovery
slot, the positive-time semiflow factors through
\begin{equation}
 Y=C([-r,0],\R)\times\R,\qquad
 \Pi_Y(\phi_v,\phi_w)=(\phi_v,\phi_w(0)).               \label{eq:reduced-space}
\end{equation}
This reduction preserves every nonzero Floquet multiplier and pulls stable
sets back exactly; an obsolete recovery history is not an additional
crossing coordinate.

\subsection{Finite row-mass Dobrushin networks}

For a finite network let \(Q,B_0,B_1\in\R^{N\times N}\) be nonnegative
matrices and let \(\pi\) be any stationary probability vector satisfying
\begin{equation}
 Q\mathbf1=\mathbf1,\quad \pi^TQ=\pi^T,\quad
 \pi^T\mathbf1=1,\quad
 \tau(Q)\le\frac12,\qquad
 B_j\mathbf1=\frac12\mathbf1.                           \label{eq:row-mass-network}
\end{equation}
Here \(\tau(Q)\) is the Dobrushin coefficient.  Zero entries of \(\pi\)
are allowed, and no lower bound on \(\pi_{\min}\) enters the estimates.
The row-mass identities preserve the collective line
\(\Span\{\mathbf1\}\), but \(\Ker\pi^T\) need not be invariant under the
delayed layers: no identity \(\pi^TB_j=\frac12\pi^T\) is assumed.
Collective/quotient coordinates therefore produce an upper-triangular
operator.  The former balanced direct sum is the special case
\[
 \delta_B=0,\qquad
 \delta_B:=\delta_0+\delta_1,\qquad
 \delta_j:=\frac12\norm{\pi^TB_j-\tfrac12\pi^T}_1,
 \qquad 0\le\delta_B\le1.
\]

The finite leaky network to which the estimates below apply is, with
\(f(s)=s-s^3/3\) and both \(f\) and \(H\) acting componentwise,
\begin{align}
 \dot v(t)={}&f(v(t))-w(t)+3(Q-\Id)v(t)\notag\\
 &+\eps\kappa_1\{B_0v(t-\tau_0)+B_1v(t-\tau_1)-v(t)\}\notag\\
 &+\eps\kappa_3\{B_0H(v(t-\tau_0))+B_1H(v(t-\tau_1))-H(v(t))\},
                                                               \label{eq:leaky-network-v}\\
 \dot w(t)={}&\eps\{v(t)-a\mathbf1-w(t)\}+2(Q-\Id)w(t).
                                                               \label{eq:leaky-network-w}
\end{align}
Here \(v,w\in\R^N\), and all scalar constants are exactly those in
\cref{eq:center}.  A collective physical pulse adds
\(J\mathbf1_{[0,1)}(t)\mathbf1\) to
\cref{eq:leaky-network-v}.  The row-mass identities imply that
\(v=V\mathbf1\), \(w=W\mathbf1\) is invariant and that its restriction is
exactly \cref{eq:scalar-v,eq:scalar-w}, including the two half-weighted
delayed terms.  Thus the constants in the transverse theorem below refer to
this RFDE, not to an unspecified network realization.

Use the transverse history diameter norm
\begin{equation}
 \norm{\phi}_{1/10}
 =\sup_{-r\le\theta\le0}e^{\theta/10}
   \max\{\osc\phi_v(\theta),3\osc\phi_w(\theta)\}.      \label{eq:diameter-norm}
\end{equation}

\begin{theorem}[Complete-line transverse inverse]\label{thm:transverse}
Along every complete synchronous leaky trajectory satisfying
\(\sup_t\abs{V(t)-1}\le5/2\), every network satisfying
\cref{eq:row-mass-network} has
\begin{equation}
 \norm{U_{\perp,N}(t,s)}_{1/10}\le e^{-(t-s)/10},
 \qquad \norm{G_{\perp,N}}\le10.                       \label{eq:green-bound}
\end{equation}
The constants are independent of the finite network size and admitted
topology.
\end{theorem}

The proof uses a directed Dobrushin--Halanay estimate and a forward pullback
of zero-history forced solutions.  It never inverts an RFDE semiflow on the
whole history space.

\begin{corollary}[Upper-triangular canonical root transfer]
\label{cor:canonical-transfer}
Suppose a scalar phase-fixed Lin operator has Fredholm index \(-1\), a
one-dimensional normalized cokernel, gap \(d(\nu)\), and a simple root
\(d(\nu_c)=0\).  If all network endpoint, phase, parameter, and gap rows are
collective, then
\begin{equation}
 \cL_N=
 \begin{pmatrix}\cL_{\parallel}&C_N\\0&\cL_{\perp,N}\end{pmatrix},
 \qquad
 T_N\cL_N=\cL_{\parallel}\oplus\cL_{\perp,N},           \label{eq:lin-sum}
\end{equation}
where
\[
 T_N(y_\parallel,y_\perp)
 =\bigl(y_\parallel-C_NG_{\perp,N}y_\perp,y_\perp\bigr),
 \qquad
 \norm{C_NG_{\perp,N}}\le\frac{391}{2000}\delta_B.
\]
If \(\psi\) is the normalized scalar cokernel, the full normalized
cokernel is
\begin{equation}
 \Psi_N(y_\parallel,y_\perp)
 =\psi\bigl(y_\parallel-C_NG_{\perp,N}y_\perp\bigr).    \label{eq:full-cokernel}
\end{equation}
For a collective parameter row \((g_\parallel,0)\), one has
\(\Psi_N(g_\parallel,0)=\psi(g_\parallel)\).  Hence
\(d_N(\nu)=d(\nu)\), and every admitted finite network has the same
Fredholm index, root, slope, and orientation as the scalar problem.
\end{corollary}

For the leaky model, \cref{thm:transverse} verifies the network block of
this corollary.  The result is conditional on the scalar phase-fixed
complete-history root and the declared collective preparation.  The scalar
leaky root itself, an invariant asynchronous voltage strip, and physical
pulse onset remain separate open problems.

\subsection{Uniform structural transfer}

Let \(\widetilde d_N(\nu,\cR)\) be a normalized selected gap on a common
parameter cylinder.  The residual norm is an operator norm controlling the
instantaneous block, fixed-support delay measure, nonlinear jets, and trace
preparation; it is not an unscaled Frobenius norm.

\begin{theorem}[Dimension-uniform one-gap theorem]\label{thm:one-gap}
Assume the graph--trace or augmented-Lin range problem is uniformly solvable,
the model-fitting error is at most \(C_{\rm fit}\norm{\cR}\), and
\(\widetilde d_N\) is jointly \(C^2\).  Suppose, uniformly in \(N\),
\begin{equation}
 \widetilde d_N(\nu_{0,N},0)=0,\qquad
 \abs{a_N}:=\abs{\partial_\nu\widetilde d_N(\nu_{0,N},0)}\ge m_*>0,
 \quad \norm{D_{\cR}\widetilde d_N}\le B_*,
 \quad \norm{D^2\widetilde d_N}\le M_*.               \label{eq:root-bounds}
\end{equation}
Here \(B_*,M_*,C_{\rm fit}\ge0\).  If
\(r_\nu,r_{\cR},\eta_*>0\) are respectively the common cylinder radii and
the range-solve radius, set
\begin{equation}
 \rho_*=\min\left\{r_{\cR},\frac{m_*r_\nu}{2B_*},
 \frac{m_*}{2M_*(1+2B_*/m_*)},
 \frac{\eta_*}{C_{\rm fit}}\right\}.                   \label{eq:rho-star}
\end{equation}
Every quotient in \cref{eq:rho-star} whose nonnegative denominator is zero
is, by convention, \(+\infty\); hence an identically absent error term
imposes no artificial radius restriction.
For every \(r=\norm{\cR_N}\le\rho_*\), there is exactly one root in
\begin{equation}
 \abs{\nu-\nu_{0,N}}\le\frac{2B_*}{m_*}r,              \label{eq:root-ball}
\end{equation}
and, with \(\ell_N=D_{\cR}\widetilde d_N(\nu_{0,N},0)\),
\begin{equation}
 \left|\nu_{c,N}(\cR_N)-\nu_{0,N}
       +\frac{\ell_N[\cR_N]}{a_N}\right|
 \le\frac{M_*}{2m_*}\left(1+\frac{2B_*}{m_*}\right)^2r^2. \label{eq:root-remainder}
\end{equation}
All constants are independent of the admitted finite network size.
\end{theorem}

The theorem concerns a selected complete-history root.  A zero-fiber
implication is additionally required to call it a geometric canard, and a
basin theorem is required to call anything a physical pulse threshold.

\subsection{A synchrony-quotient-free instance}

We next state a separate network result whose node fields are heterogeneous
and whose slow recovery variable is shared.  For each \(N\ge2\), let
\(P_N\ge0\) be row stochastic and let \(\pi_N\) be a strictly positive
stationary probability vector.  For fixed \(D,\gamma>0\), assume
\begin{equation}
 P_N\mathbf1=\mathbf1,\qquad
 \pi_N^TP_N=\pi_N^T,\qquad \pi_N^T\mathbf1=1,\qquad
 \tau(P_N)\le1-\gamma.                                \label{eq:shared-markov}
\end{equation}
Set
\begin{equation}
 P_{c,N}=\mathbf1\pi_N^T,\qquad P_{\perp,N}=\Id-P_{c,N},\qquad
 E_N=\Ker\pi_N^T,\qquad A_N=D(P_N-\Id)|_{E_N},       \label{eq:shared-splitting}
\end{equation}
and equip \(\R^N\) with
\begin{equation}
 \norm{x}_N=\abs{\pi_N^Tx}+\osc(x).                   \label{eq:shared-norm}
\end{equation}
Let the curvature vectors satisfy, with constants independent of \(N\),
\begin{equation}
 0<c_-\le c_{i,N}\le c_+,\qquad \pi_N^Tc_N=\alpha>0. \label{eq:shared-curvature}
\end{equation}
Fix \(\beta>0\), \(K\ne0\), and finitely many delay locations
\(0\le\theta_0<\cdots<\theta_m\le\Theta_*\).  For
\begin{equation}
 B_{k,N}(\zeta)=B_{k,N}+\zeta R_{k,N},\qquad
 C_N(\zeta)=\sum_{k=0}^mB_{k,N}(\zeta),                \label{eq:shared-layers}
\end{equation}
assume the fixed-support operator bound and the atomwise full-row
neutrality
\begin{equation}
 \sup_N\sum_{k=0}^m
 \left(\norm{B_{k,N}}_{\cL(\R^N,\norm{\cdot}_N)}
       +\norm{R_{k,N}}_{\cL(\R^N,\norm{\cdot}_N)}\right)<\infty,
 \qquad \pi_N^TR_{k,N}=0.                             \label{eq:shared-layer-bounds}
\end{equation}
The latter is a row identity, not merely
\(\pi_N^TR_{k,N}\mathbf1=0\).

For \(\eps=\delta^2\), the RFDE is
\begin{align}
 \dot v(t)={}&\frac23\mathbf1-w(t)\mathbf1
 -c_N\circ(v(t)-\mathbf1)^{\circ2}
 -\frac\beta3(v(t)-\mathbf1)^{\circ3}+D(P_N-\Id)v(t)\notag\\
 &+\delta^2K\left\{C_N(\zeta)v(t)
 -\sum_{k=0}^mB_{k,N}(\zeta)
       v\left(t-\frac{\theta_k}{\delta}\right)\right\},
                                                               \label{eq:shared-rfde-v}\\
 \dot w(t)={}&\delta^2\{\pi_N^Tv(t)-a\},             \label{eq:shared-rfde-w}
\end{align}
where \(v\in\R^N\), \(w\in\R\), and the powers and products are
componentwise.  Its collective fold is
\((v,w,a)=(\mathbf1,2/3,1)\).  Write
\begin{equation}
 a=1+\delta^2\nu,\qquad \mu=a-1=\delta^2\nu,\qquad
 \dot M_{1,N}=\sum_{k=0}^m\theta_kR_{k,N},             \label{eq:shared-parameters}
\end{equation}
and define
\begin{align}
 g_N&=A_N^{-1}P_{\perp,N}c_N,\notag\\
 \kappa_N&=\frac\beta3+2\pi_N^T\diag(c_N)g_N,\qquad
 \nu_{0,N}=-\frac{3\kappa_N}{8\alpha^3},             \label{eq:shared-center}\\
 r_N&=\frac{K}{2\alpha}\pi_N^T\diag(c_N)
 A_N^{-1}P_{\perp,N}\dot M_{1,N}\mathbf1.            \label{eq:shared-rn}
\end{align}

For completeness, we also fix the selection entering the theorem.  In the
fold chart
\[
 s=\delta t,\qquad w=\frac23-\delta^2Y,\qquad
 v=\mathbf1+\delta\mathbf1X
       +\delta^2\{g_NX^2+h\},\qquad h\in E_N,
\]
the singular critical field is
\(q_0(X,Y)=(Y-\alpha X^2,-X)\).  A uniformly admissible preparation rule
\(\mathcal P=(\mathcal P_N)\) means that the same frozen graph cutoffs,
planar joining cutoff, degree-three normal extension, phase buffer, and
invariant-tail levels are used for every \(N\), with common derivative
bounds; the preparation is independent of \((\nu,\zeta)\) and agrees with
\cref{eq:shared-rfde-v,eq:shared-rfde-w} on the retained continuous
depth-two \(q_0\)-flow hull of half-length
\(S_\delta=\sqrt{2p\log(1/\delta)}\), for an exponent \(p>4\) fixed in the
theorem below.  It selects attracting and repelling phase-fixed planar
traces \(z^a_{N,\mathcal P_N}\) and
\(z^r_{N,\mathcal P_N}\), with phase \(X(0)=0\).  Their normalized central
gap is
\begin{align}
 \mathscr D_{N,\mathcal P_N}(\delta,\nu,\zeta)
  &=\frac{2}{\alpha e}\left\{
    \mathscr H_\alpha(z^a_{N,\mathcal P_N}(0))
   -\mathscr H_\alpha(z^r_{N,\mathcal P_N}(0))\right\},
                                                               \label{eq:shared-gap}\\
 \mathscr H_\alpha(X,Y)
  &=\frac12e^{-2\alpha Y}
    \left(\alpha Y-\alpha^2X^2+\frac12\right).       \label{eq:shared-first-integral}
\end{align}
On the retained graph, a zero of \(\mathscr D_{N,\mathcal P_N}\) is
equivalent to equality of the two selected complete RFDE histories.

\begin{theorem}[Shared-resource network response]\label{thm:shared-resource}
Assume \cref{eq:shared-markov,eq:shared-curvature,eq:shared-layer-bounds}
with all displayed bounds common in \(N\), and assume
\begin{equation}
 \inf_{N\ge2}\abs{r_N}>0.                              \label{eq:shared-nondegeneracy}
\end{equation}
Fix \(p>4\) sufficiently large, a compact interval \(I_\nu\) containing
every \(\nu_{0,N}\) in its interior with one common positive boundary
distance, and a class of uniformly admissible preparation rules with one
set of bounds.  Then there exist
\(\delta_0,\zeta_0,c_0,C>0\), independent of \(N\) and of the rule in
that class, such that for every \(N\ge2\), every
\(0<\delta<\delta_0\), every \(\abs\zeta<\zeta_0\), and every admitted
\(\mathcal P\), the gap \(\mathscr D_{N,\mathcal P_N}/\delta\) has exactly
one root in the local window
\(\{\nu:\abs{\nu-\nu_{0,N}}<c_0\}\subset I_\nu\),
\begin{equation}
 \nu_{c,N,\mathcal P_N}(\delta,\zeta)\in I_\nu,
 \qquad
 \abs{\nu_{c,N,\mathcal P_N}(\delta,\zeta)-\nu_{0,N}}<c_0. \label{eq:shared-root}
\end{equation}
No assertion is made here about additional zeros elsewhere in \(I_\nu\).
The selected attracting and repelling complete histories agree at this
root.  With
\(\mu_{c,N,\mathcal P_N}=\delta^2\nu_{c,N,\mathcal P_N}\),
\begin{align}
 &\left|\mu_{c,N,\mathcal P_N}(\delta,\zeta)
       -\mu_{c,N,\mathcal P_N}(\delta,0)
       -\mathscr C_N\delta^3\zeta\right|
 \le C\{\delta^4\abs\zeta+\delta^3\zeta^2\},        \label{eq:shared-response}\\
 &\mathscr C_N=-\frac{r_N}{\alpha}
 =-\frac{K}{2\alpha^2}\pi_N^T\diag(c_N)
 A_N^{-1}P_{\perp,N}\dot M_{1,N}\mathbf1.             \label{eq:shared-coeff}
\end{align}
In particular,
\(\inf_N\abs{\mathscr C_N}>0\).  The exact finite-\(\delta\) root may
depend on \(\mathcal P_N\); its coefficient and the displayed remainder
bound do not within the declared preparation class.
\end{theorem}

The nondegeneracy class is nonempty without a hidden synchrony quotient.
Indeed, take \(\pi_N=N^{-1}\mathbf1\), \(P_N=P_{c,N}\), and
\[
 s_{i,N}=\frac{\sqrt3(2i-N-1)}{\sqrt{(N-1)(N+1)}},
 \qquad c_N=\mathbf1+\sigma s_N,\qquad 0<\sigma<1/\sqrt3.
\]
These choices satisfy
\[
 \pi_N^Ts_N=0,\qquad \pi_N^Ts_N^{\circ2}=1,\qquad
 \norm{s_N}_\infty<\sqrt3,\qquad \alpha=1.
\]
For two delays \(\theta_0<\theta_1\), set
\[
 B_{0,N}=B_{1,N}=bP_{c,N},\qquad
 R_{0,N}=s_N\pi_N^T,\qquad R_{1,N}=-s_N\pi_N^T,
 \qquad b>0.
\]
The layers remain entrywise positive for
\(\abs\zeta<b/\sqrt3\), all entries of
\(c_N\) are distinct, and
\begin{equation}
 \mathscr C_N=\frac{K\sigma(\theta_0-\theta_1)}{2D}\ne0
                                                               \label{eq:shared-witness-coeff}
\end{equation}
for every \(N\ge2\).  The distinct curvatures destroy every nontrivial
synchrony polydiagonal.

This instance proves that the general root mechanism is not merely an
arbitrary-size replication of a two-module quotient.  It does not cover
moving delay atoms, arbitrary perturbations of the base topology, or graph
families whose mixing gap closes, and it does not identify the
preparation-indexed root with an independently selected physical outer
canard.

\subsection{Eventually smooth selected returns}

The next result is independent of the preceding network reduction.  It
isolates the regularity needed to turn a selected event into a moving
\emph{complete-history} return.  This distinction matters because a
fixed-time RFDE solution map may be smooth in its initial history while its
joint dependence on time and an arbitrary continuous history is not yet
smooth: before the translation part has cleared the initial interval,
changing time moves an evaluation point in a merely continuous function.

Let
\begin{equation}
 \mathcal X_{\tau_*}=C([-\tau_*,0],\R^d),\qquad
 \dot x(t)=\mathcal F(x_t),
 \qquad x_t(\theta)=x(t+\theta),                      \label{eq:general-rfde}
\end{equation}
where \(d<\infty\), \(\tau_*>0\), and
\(\mathcal F\colon U\subset\mathcal X_{\tau_*}\to\R^d\) is \(C^r\).  This includes any
finite network with finitely many constant delays bounded by \(\tau_*\).
Write \(\Phi_t(\phi)=x_t(\phi)\).

\begin{theorem}[Eventually smooth selected-event return]
\label{thm:eventual-selected-return}
Let \(1\le k\le r\), let \(\mathcal M\) be a Banach space,
\(D\subset\mathcal M\) open, and let
\(\iota\colon D\to U\subset\mathcal X_{\tau_*}\) be \(C^k\).  Suppose all
solutions from \(\iota(D)\) exist in one common \(U\)-tube through \(T_+\).
Let \(V\subset\mathcal X_{\tau_*}\) be open, let
\(g\colon V\to\R\) be \(C^k\), and assume
\(\Phi_t(\iota(u))\in V\) for \((t,u)\in I\times D\).  Put
\[
 \mathcal H(t,u)=g(\Phi_t(\iota(u))).
\]
Suppose a common interval \(I=[T_-,T_+]\) satisfies
\begin{equation}
 T_->k\tau_*,\qquad
 \sup_{u\in D}\mathcal H(T_-,u)\le-\delta_-,\qquad
 \inf_{u\in D}\mathcal H(T_+,u)\ge\delta_+,
 \qquad
 \partial_t\mathcal H\ge a_*>0                     \label{eq:eventual-return-gates}
\end{equation}
on \(I\times D\), for \(\delta_-,\delta_+,a_*>0\).  Then each
\(u\in D\) has exactly one selected event
\(T(u)\in(T_-,T_+)\), and the event time and ambient event-hit map
\[
 T\colon D\to\R,
 \qquad \mathcal E(u)=\Phi_{T(u)}(\iota(u))
       \colon D\to\mathcal X_{\tau_*}
\]
are \(C^k\).  The threshold \(k\tau_*\) is a uniform sufficient condition;
no necessity or optimality is asserted.  It is independent of the number of
nodes in any finite network, although concrete tube and derivative constants
need not be.

If \(\iota\) is a \(C^k\) chart of a local section \(\Sigma\), and
\(\mathcal E(D_0)\subset\iota(D)\) for an open \(D_0\subset D\), then
\begin{equation}
 \mathcal Q=\iota^{-1}\circ\mathcal E\big|_{D_0}
 \colon D_0\to D                                      \label{eq:induced-selected-return}
\end{equation}
is the induced \(C^k\) selected section-return map.
\end{theorem}

\begin{proof}
On the common tube, repeated Fr\'echet differentiation of the Volterra
equation gives, for every fixed time, the triangular initial-data
variational hierarchy through order \(k\).  We record the additional joint
smoothing step because it is essential here.  Inductively in \(j\), the
mixed jets
\[
 \partial_t^aD_\phi^b x(t;\phi),
 \qquad a+b\le j,
\]
exist after \((j-1)\tau_*\) and depend continuously on \((t,\phi)\) in the
corresponding multilinear-operator norms.  For \(j=1\) this follows from the
integral equation and the first variational equation.  If it holds through
order \(j\), differentiating \(x'=\mathcal F(x_t)\) and the triangular
variational equations once more expresses every order-\(j+1\) jet as a
finite sum of continuous derivatives of \(\mathcal F\) applied to lower-order
history jets.  After one further maximum delay, every argument segment lies
in the already regularized interval.  The linear inhomogeneous variational
equations and Gronwall's inequality give continuity in the multilinear-
operator norms, completing the induction.

The earliest physical time in \(\Phi_t(\phi)\) is \(t-\tau_*\).  Hence, for
every total order \(j\le k\), the preceding jets define continuous
\(\mathcal X_{\tau_*}\)-valued derivatives once \(t>j\tau_*\).  Therefore
\((t,\phi)\mapsto\Phi_t(\phi)\) is jointly \(C^k\) for \(t>k\tau_*\), and
\(\mathcal H\) is jointly \(C^k\) on \(I\times D\).  The endpoint signs and
positive speed give existence and uniqueness in the declared window; the
Banach implicit-function theorem gives \(T\in C^k\), and composition gives
\(\mathcal E\in C^k\).  The last assertion follows by composing with the
section chart inverse.  See also the fixed-time RFDE and Poincar\'e-map
background in \citet[Chs.~VII.4 and XIV.3]{diekmann1995delay}.
\end{proof}

No exclusion of earlier hits is used in
\cref{thm:eventual-selected-return}: it is needed only to call the branch a
first return or to assign it an exact ordinal.  This separation also avoids
an unnecessary algebraic restriction on a stable-manifold construction.

\begin{lemma}[Direct selected return identifies the stable-set germ]
\label{lem:direct-selected-stable-germ}
Let \(\Phi\) be a continuous semiflow, \(\Gamma\) a periodic orbit through a
transverse local section \(\Sigma\) at \(p\).  Choose a local section patch
\(N\subset\Sigma\) with the phase-isolation property
\begin{equation}
 x_n\in N,\qquad \operatorname{dist}(x_n,\Gamma)\to0
 \quad\Longrightarrow\quad x_n\to p.                 \label{eq:section-phase-isolation}
\end{equation}
Let \(\mathcal Q(x)=\Phi_{\Theta(x)}(x)\) be a selected local return of the
same semiflow with
\[
 \mathcal Q\colon N\to N,\qquad \mathcal Q(p)=p,\qquad
 \Theta(p)=mP,\qquad
 0<\Theta_-\le\Theta(x)\le\Theta_+<\infty.
\]
Suppose every intervening arc between successive iterates stays in one common
flow tube.  Restrict both stable sets to points whose relevant iterates or
flow arcs stay in these declared local domains, and assume that local
flow-stable points in \(N\) have the selected hits recurrently defined in
\(N\).  Then, as germs at \(p\),
\begin{equation}
 \operatorname{germ}_p W^s_N(\mathcal Q)
 =\operatorname{germ}_p\bigl(N\cap W^s_{\rm tube}(\Gamma)\bigr).
                                                               \label{eq:direct-return-stable-germ}
\end{equation}
Consequently, a direct near-\(mP\) selected return can construct the intrinsic
periodic-orbit stable-sheet germ; an identity \(\mathcal Q=P^m\) and a
first-return label are sufficient conveniences, not necessary hypotheses.
If \(\mathcal Q\) is \(C^2\), \(p\) has the required hyperbolic splitting,
and a local graph theorem applies, its stable graph represents this same
intrinsic germ.
\end{lemma}

\begin{proof}
Put
\(t_n=\sum_{j=0}^{n-1}\Theta(\mathcal Q^j x)\).  If
\(\mathcal Q^n(x)\to p\), then \(t_n\to\infty\) by the positive lower bound;
joint continuity of \(\Phi_t\), uniformly for
\(0\le t\le\Theta_+\), sends every intervening late-time arc to the
corresponding arc of \(\Gamma\).  Conversely, convergence of the flow orbit
to \(\Gamma\), together with recurrent selected hits in \(N\), gives
\(\operatorname{dist}(\mathcal Q^n x,\Gamma)\to0\).  The phase-isolation
property \cref{eq:section-phase-isolation} then gives
\(\mathcal Q^n x\to p\).
\end{proof}

The theorem deliberately excludes state-dependent or neutral delays,
infinite delays, infinite networks, and infinite-dimensional node states.
It proves a reusable return-map mechanism, not a return tube or event window
for the concrete leaky model.

\subsection{Conditional physical-pulse and control theorem}

Let \(\xi=(a,\kappa_3)\) range in a closed rectangle
\(\Xi\subset\R^2\) about
\(\xi_0=(1/4,1/200)\), with \(\eps,\kappa_1,\tau_0,\tau_1\) fixed as in
\cref{eq:center}.  Write
\begin{equation}
 \bar v_q(a)=(3a)^{1/3},\qquad
 E_q(\xi)(\theta)\equiv
 \bigl(\bar v_q(a),\bar v_q(a)-a\bigr),
 \quad -r\le\theta\le0,                               \label{eq:quiet-equilibrium}
\end{equation}
for the quiet equilibrium history in \(X\), and let
\(\Phi_t^\xi\) be the autonomous post-release semiflow.  Let
\(\Gamma_i(\xi)\) and \(\Gamma_o(\xi)\) be the inner and outer scalar
periodic orbits, identified with their compact orbits of histories in
\(X\), and set
\begin{equation}
 F(\xi)=T_o(\xi)^{-1},\qquad
 A(\xi)=\max_tV_o(t;\xi)-\min_tV_o(t;\xi).             \label{eq:outer-observables}
\end{equation}
For a physical one-unit pulse of amplitude \(J\), applied from
\(E_q(\xi)\), denote the history at release by \(R_\xi(J)\in X\).  On a
phase section \(\Sigma_\xi\subset Y\) through
\(\Gamma_i(\xi)\), let \(B_\xi(J)\) be the event-aligned entrance of the
post-release orbit.  In section coordinates
\((u,z)\in E^u_\xi\oplus E^s_\xi\), write the local stable sheet as
\begin{equation}
 W^s_{\rm loc}(\Gamma_i(\xi))\cap\Sigma_\xi
   =\{(u,z):u=\psi_\xi(z)\},
 \qquad H(\xi,J)=u_\xi(J)-\psi_\xi(z_\xi(J)),          \label{eq:pulse-gap}
\end{equation}
where \((u_\xi(J),z_\xi(J))\) are the coordinates of \(B_\xi(J)\).

\begin{theorem}[Pulse threshold and safety coordinates: completion
theorem]\label{thm:pulse-completion}
Assume the following objects and estimates hold with common domains, radii,
and strict constants for every \(\xi\in\Xi\).  The equilibrium
\(E_q(\xi)\) and the phase-fixed branches
\(\Gamma_i(\xi),\Gamma_o(\xi)\) are \(C^1\) in \(\xi\); the quiet
equilibrium and outer cycle possess quantitative attraction tubes; and the
outer voltage extrema are simple.  The inner section return has a
one-dimensional unstable bundle and a stable complement, both \(C^1\) in
\(\xi\), and \(\psi_\xi(z)\) in \cref{eq:pulse-gap} is jointly \(C^1\)
in \((\xi,z)\) on one common graph domain.

Assume also that \((\xi,J)\mapsto R_\xi(J)\) is \(C^1\) on
\(\Xi\times I_J\), where \(I_J=[J_-,J_+]\), and that
\(B_\xi(J)\) is one jointly \(C^1\), event-aligned entrance branch on this
whole product: the section event does not switch branches, its speed is
bounded away from zero, and every \(B_\xi(J)\) lies in the common domain of
the signed-exit construction, while every stable coordinate \(z_\xi(J)\)
lies in the common domain of \(\psi_\xi\).  Suppose, after fixing
the orientation, that there are \(\eta_H,m_J>0\) such that
\begin{equation}
 H(\xi,J_-)\ge\eta_H,\qquad H(\xi,J_+)\le-\eta_H,
 \qquad \partial_JH(\xi,J)\le-m_J<0,                 \label{eq:gap-crossing}
\end{equation}
for all \(\xi\in\Xi\) and \(J\in I_J\).  In particular, \(H\) is a
jointly \(C^1\) function on this full product, rather than a collection of
pointwise gaps.  Finally, assume the signed-exit
dynamics sends the \(H>0\) side into a complete-history exit slab attached
to the quiet attraction tube and the \(H<0\) side into a complete-history
exit slab attached to the outer attraction tube, uniformly in \(\xi\),
with no other local exit; histories on \(H=0\) lie on
\(W^s_{\rm loc}(\Gamma_i(\xi))\).

Then there is a unique \(C^1\) function
\(J_c:\Xi\to(J_-,J_+)\) satisfying \(H(\xi,J_c(\xi))=0\), with
\begin{equation}
 D_\xi J_c(\xi)
 =-\frac{D_\xi H(\xi,J_c(\xi))}
          {\partial_JH(\xi,J_c(\xi))}.                 \label{eq:threshold-derivative}
\end{equation}
For every \((\xi,J)\in\Xi\times I_J\), the routing assertions mean the
following asymptotic statements in the complete history space:
\begin{equation}
 \begin{aligned}
 J<J_c(\xi)&\ \Longrightarrow\
   \norm{\Phi_t^\xi R_\xi(J)-E_q(\xi)}_X\longrightarrow0,\\
 J=J_c(\xi)&\ \Longrightarrow\
   \operatorname{dist}_X(\Phi_t^\xi R_\xi(J),\Gamma_i(\xi))
      \longrightarrow0,\\
 J>J_c(\xi)&\ \Longrightarrow\
   \operatorname{dist}_X(\Phi_t^\xi R_\xi(J),\Gamma_o(\xi))
      \longrightarrow0
 \end{aligned}
 \qquad(t\to\infty).                                  \label{eq:routing}
\end{equation}
Define the signed intrinsic safety margin
\(S(\xi,J)=J-J_c(\xi)\).  If
\(\xi_*\in\operatorname{int}\Xi\) satisfies
\(\det D_\xi(F,A)(\xi_*)\ne0\), then, for every
\(J_*\in\operatorname{int}I_J\),
\begin{equation}
 \cQ(\xi,J)=(F(\xi),A(\xi),S(\xi,J))                 \label{eq:safety-map}
\end{equation}
is a local \(C^1\) diffeomorphism at \((\xi_*,J_*)\), and
\begin{equation}
 \det D\cQ(\xi_*,J_*)=\det D_\xi(F,A)(\xi_*).        \label{eq:safety-determinant}
\end{equation}
\end{theorem}

The implication in \cref{thm:pulse-completion} is proved by a forward
Lyapunov--Perron stable graph, graph-adjusted stable-gap endpoint signs plus
strict gap monotonicity,
two-sided routing, and the block-triangular inverse theorem.  Interval Newton
may sharpen the resulting scalar root enclosure but is not needed for its
existence or uniqueness.  The model-specific graph and routing premises are
current release gates.  What is already proved
at the center is the exact reduced-history factorization, strict quiet
capture at \(J=0.30\), exactly one nontranslation unstable inner multiplier,
an explicit stable spectral gap, a continuous Route--C adjoint history
measure with nonzero normalization, a physical-time base-orbit correlated
second-return bound \(C^{uu}_{s,\mathrm{base}}<7.905650<12\), directed outer \((F,A)\)
nonsingularity, a rigorous fixed-time fourth-order pulse-parameter model
with fifth-order remainder on \(J\in[0.30105,0.30120]\), and an ambient
complete-history proximity bound at \(J=0.32\), not a section/domain
attachment.  On that pulse interval a
separate validated event theorem also gives a unique positive Route--C
crossing in a common bracket, a fourth-order event-time graph, and a
continuous common-event history tube; see \cref{thm:route-c-event}.  Its
event-aligned complete-history first derivative, including the moving-time
term, is also enclosed on the full interval in
\cref{thm:event-first-derivative}.  After correcting the complex right
eigenvector gauge, \cref{thm:oriented-pulse-action} proves that the fixed
normalized Route--C functional has a strictly negative action on this
derivative throughout the interval.  In the same continuous-history norm,
Route--C section, physical Grushin normalization, and centered coordinate
\(\kappa(J)=K(J)-X_*\), source-bound endpoint calculations prove
\(f_{\rm phys}(\kappa(J_-))>0>f_{\rm phys}(\kappa(J_+))\).  The direct
complete-history norm of the same \(q_{\rm phys}\) gives
\(\|P_sD_J\kappa\|_Y\le5.979324\), and a two-ended cone proves
\(P_s\kappa(I_J)\subset\overline B_Y(0,47/5000)\cap\ker f_{\rm phys}\).
Conditional on a quantitative graph whose stable domain contains this ball
and \(\sup\|D\psi\|\le16\), it gives
\(H'(J)\in[-354.415700,-149.584300]\subset(-\infty,0)\) on the full pulse
interval.  The general return theorem
\cref{thm:eventual-selected-return} identifies the missing \(C^2\) adapter:
the one-period center does not clear two maximum delays, whereas the
near-two-period center does so with margin greater than \(14.010740\).
\Cref{prop:explicit-two-period-event-tube} now proves a common \(C^2\)
selected event and complete-history hit on the explicit arbitrary-
continuous-history scale \(\lambda_0=9\times10^{-31}\), and the orbit-aware
enlargement in \cref{prop:orbit-aware-two-period-event-tube} raises the
proved closed scale to \(1.2\times10^{-8}\), while
\cref{prop:direct-two-period-derivative} proves directly that its center
derivative is \(A^2\).  For that route, the conditional six-block box in
\cref{prop:two-return-graph-arithmetic} would give the sharper physical graph
budgets and
\(H'(J)\in[-260.794147,-243.205853]\), together with strict graph-adjusted
endpoint margins.  Its signed finite-section pilot lies comfortably inside
all six caps, but is not interval evidence.  Thus the functional endpoint
signs, pulse ball containment, enlarged but still sub-preferred two-period
event tube, center
derivative, fixed-splitting/rate transfer, and conditional graph/crossing
arithmetic are proved.  Enlargement of the return to the preferred graph
ball, graph-transform domain containment, six directed continuous-history
Hessian bounds, stable-set identification, the graph, selected crossing,
onset, and routing remain open.  Interval Newton is only an optional later
sharpening.
Neither that observable event and derivative nor the preceding
invariant-object facts
imply \cref{eq:gap-crossing} or \cref{eq:routing}.
